\clearpage\begingroup\SourceLeaf{69b}\HistoricalHead{64}{L}{PHYSIK. NACHLASS.}
\begin{center}[7.]\end{center}
\noindent{\small [S. 214]}
\begin{center}\textbf{Inducirte gemischte Bewegung.}\end{center}
\noindent\hspace*{\parindent}\begin{tabular}{@{}l@{\quad}l@{}}
Schwingungsdauer des inducirenden Stabes&\(\dfrac{180^0}{n}=60^{\prime\prime}4\)\\[1ex]
Schwingungsdauer des passiven Stabes&\(\dfrac{180}{N}=20^{\prime\prime}6\)
\end{tabular}\par
\indent Dadurch wird die gemischte Bewegung inducirt, wenn die Verbindung
zur Zeit \(\theta\) anfängt, und für die Zeit \(0\) der inducirende Stab mit positiver
Bewegung durch die Mitte geht:
\[
\begin{aligned}
&(N\cos n\theta\sin N\theta-n\sin n\theta\cos N\theta)\sin Nt\\
&+(N\cos n\theta\cos N\theta+n\sin n\theta\sin N\theta)\cos Nt\\
&-N\cos nt.
\end{aligned}
\]
\noindent Hebt man die Verbindung zur Zeit \(\theta'\) auf und schliesst sie wieder zur Zeit
\(\theta''\), so ist die gemischte Bewegung hernach
\clearpage\SourceLeaf{70}\HistoricalHead{65}{R}{BERECHNUNG MAGNETISCHER KRÄFTE.}\HistoricalFoot{XI1.}{9}
\[
\begin{aligned}
&\left\{\begin{aligned}
&N\cos n\theta\sin N\theta-n\sin n\theta\cos N\theta\\
&-N\cos n\theta'\sin N\theta'+n\sin n\theta'\cos N\theta'\\
&+N\cos n\theta''\sin N\theta''-n\sin n\theta''\cos N\theta''
\end{aligned}\right\}\sin Nt\\[1ex]
&+\left\{\begin{aligned}
&N\cos n\theta\cos N\theta+n\sin n\theta\sin N\theta\\
&-N\cos n\theta'\cos N\theta'-n\sin n\theta'\sin N\theta'\\
&+N\cos n\theta''\cos N\theta''+n\sin n\theta''\sin N\theta''
\end{aligned}\right\}\cos Nt\\[1ex]
&-N\cos nt.
\end{aligned}
\]
\noindent Um die Coefficienten von \(\sin Nt\) und \(\cos Nt\) verschwinden zu machen, ist
am einfachsten
\[
\begin{aligned}
\theta&=-u,\\
\theta'&=0,\\
\theta''&=+u
\end{aligned}
\]
\noindent zu setzen, wodurch der ersten Bedingung von selbst Genüge geschieht, und
die zweite noch erfordert:
\[
N(2\cos nu\mathpunct{.}\cos Nu-1)+n\mathpunct{.}2\sin nu\mathpunct{.}\sin Nu=0
\]
\noindent oder
\[
(N+n)\cos(N-n)u+(N-n)\cos(N+n)u=N.
\]
\noindent Für unsere Zahlen ergiebt sich
\[
\begin{alignedat}{2}
(N+n)u&=0^{\mathrm R}{,}977766\qquad&Nu&=0^{\mathrm R}{,}729095\\
(N-n)u&=0^{\phantom{\mathrm R}}{,}480425&nu&=0^{\phantom{\mathrm R}}{,}248671\\
&&u&=7^{\prime\prime}5096.
\end{alignedat}
\]
\bigskip\begin{center}\textbf{BEMERKUNGEN ZU [6.] UND [7.].}\end{center}
\begingroup\small
\indent Bezeichnet in [6.] \(x\) die Elongation einer der beiden Nadeln, \(n\) ihre Schwingungszahl in \(2\pi\) Sekunden
(Frequenz), und wird durch \(A\) eine der ablenkenden Kraft des Stromes proportionale Größe bezeichnet, so
lautet die Differentialgleichung der Schwingung offenbar:
\[
\frac{d^2x}{dt^2}+n^2x=A,
\]
\noindent und nach einer Kommutation:
\[
\frac{d^2x}{dt^2}+n^2x=-A.
\]

Genau ebenso für die andere Nadel, wo nur anstatt \(n\) der Wert \(n'\) zu setzen ist. Man hat die Lösungen
der beiden Gleichungspaare, die abwechselnd gelten, dann so aneinander zu setzen, daß zu den
Zeiten der Kommutationen Elongationen und Geschwindigkeiten stetig in einander übergehen. Dann erhält
man unmittelbar die Formeln in [6.].
\clearpage\SourceLeaf{71a}\HistoricalHead{66}{L}{PHYSIK. NACHLASS.}
\indent Ein allgemeineres Problem — nämlich das der erzwungenen Schwingungen — wird in [7.] behandelt.
Hier sind abwechselnd die Lösungen der folgenden Differentialgleichungen aneinander zu fügen:
\begin{center}
\begin{tabular}{@{}l@{\quad}p{.64\linewidth}@{}}
\(\dfrac{d^2x}{dt^2}+N^2x=0,\)&wenn die Verbindung zur induzierenden Nadel aufgehoben ist,\\[1.2ex]
\(\dfrac{d^2x}{dt^2}+N^2x=\cos nt,\)&wenn die Verbindung mit der induzierenden Nadel hergestellt wird.
\end{tabular}
\end{center}
\noindent Für die Bewegung der »passiven« Nadel erhält man dann leicht die angegebene Formel in [7.], in der die
beiden Glieder mit \(\sin Nt\) und \(\cos Nt\) die »\textls[120]{Eigenschwingung}«, das Glied mit \(\cos nt\) die »\textls[120]{erzwungene}«
Schwingung darstellt.

Beide Aufgaben stehen offenbar in enger Beziehung zur sogenannten Multiplikations- und Zurückwerfungsmethode,
die auf \textsc{Gauss} und \textsc{Weber} zurückgeht und in spezieller Form von \textsc{Weber}\textsuperscript{*}) veröffentlicht
ist.\begingroup\renewcommand{\thefootnote}{}\footnotetext{\textsuperscript{*}) \textit{Elektrodynamische Massbestimmungen, insbesondere Widerstandsmessungen}, Abhandlungen der Kgl.
Sächs. Gesellschaft d. Wissensch., math.-phys. Classe 1, Leipzig 1852, S. 199; \textsc{Wilh. Webers} Werke III,
Abhandl. X, S. 301, insbesondere S. 438 und 441.}\endgroup\ Zu bemerken ist übrigens, daß in der Gleichung für die »gemischte Bewegung« in [7.] ein für \textsc{Gauss}’
Zweck irrelevanter (»Resonanz-«)Faktor durchweg fortgelassen ist. — Man vergleiche auch \textsc{Gauss}’ Abhandlung
\textit{Über ein Mittel die Beobachtung von Ablenkungen zu erleichtern} in den Resultaten u.s.w. 1839,
S. 52, Werke V, S. 395.
\par\raggedleft\textsc{Schaefer}.\par
\clearpage\endgroup\endgroup
